% Generated by tools/build_cv.py from the career repo (cv/cv.yaml). Do not edit.

\cvsection{Open Source Software}

\begin{cventries}

  \cventry
    {Maintainer, main contributor}
    {\hrefs{https://github.com/aicenter/Ridesharing_DARP_instances}{DARP Instances}}
    {PostgreSQL, Python}
    {2023--Present}
    {
      \begin{cvitems}
        \item {Instances for large-scale ridesharing (DARP) problems together with the code for generating them from open data.}
        \item {Instances combine data from OpenStreetMap, municipality data, and Uber Movement.}
        \item {Three areas are covered: New York City, Chicago, and Washington, D.C.}
      \end{cvitems}
    }

  \cventry
    {Maintainer, main contributor}
    {\hrefs{https://github.com/aicenter/DARP-benchmark}{DARP Benchmark}}
    {\CC{}, Python}
    {2019--Present}
    {
      \begin{cvitems}
        \item {Framework for comparing the performance of DARP solvers. Can be used also for similar problems, e.g., fleet sizing.}
        \item {Ultra-fast Insertion Heuristic is implemented as a baseline.}
      \end{cvitems}
    }

  \cventry
    {Maintainer, main contributor}
    {\hrefs{https://github.com/aicenter/simod}{SiMoD}}
    {Java, Python}
    {2016--Present}
    {
      \begin{cvitems}
        \item {Simulation of Mobility-on-Demand Systems (SiMoD) is a software tool for simulating the operation of mobility-on-demand systems.}
        \item {The tool consists of an event-based simulation in Java and supporting Python scripts for input, output, and analysis.}
        \item {The simulation contains implementations of various algorithms for dispatching, rebalancing. There is also a station positioning algorithm in the Python part.}
      \end{cvitems}
    }

  \cventry
    {Maintainer, main contributor}
    {\hrefs{https://github.com/F-I-D-O/line_planning}{Line planning}}
    {Python}
    {2025--Present}
    {
      \begin{cvitems}
        \item {Pipeline for MoD-aware public transort line selection.}
      \end{cvitems}
    }

  \cventry
    {Maintainer, main contributor}
    {\hrefs{https://github.com/aicenter/road-graph-tool}{Road Graph Tool}}
    {Python, PostgreSQL}
    {2023--Present}
    {
      \begin{cvitems}
        \item {Tool for processing OpenStreetMap data and other sources into road graphs and instaces for various transportation problems.}
        \item {The tool is database first, Python is used for automation only}
      \end{cvitems}
    }

  \cventry
    {Maintainer, contributor}
    {\hrefs{https://github.com/aicenter/ShoDi}{Shortest Distances}}
    {\CC{}}
    {2019--Present}
    {
      \begin{cvitems}
        \item {Tool and library for fast computation of shortest distances in road networks.}
        \item {Various methods and datastructures are implemented with different trade-offs between memory usage and speed.}
        \item {The tool consists of a preprocesor for creating the datastructures and a query tool for computing the distances.}
        \item {There is a Java binding for the library.}
      \end{cvitems}
    }

  \cventry
    {Maintainer, contributor}
    {\hrefs{https://github.com/aicenter/agentpolis}{Agentpolis}}
    {Java, Python}
    {2016--Present}
    {
      \begin{cvitems}
        \item {Event-based transportation simulation framework.}
        \item {Main dependency of SiMoD and similar projects}
        \item {Versatile and highly configurable.}
      \end{cvitems}
    }

  \cventry
    {Maintainer, main contributor}
    {\hrefs{https://github.com/F-I-D-O/info-maven-plugin}{Info plugin for Maven}}
    {Java}
    {2024--Present}
    {
      \begin{cvitems}
        \item {Maven plugin for listing dependencies cached in the local repository.}
        \item {Available at Maven Central.}
      \end{cvitems}
    }

\end{cventries}
